\section{Purpose and summary}
\label{sec:intro}

The diagnostic under development detects protease activity in wound fluid. A sensor peptide is cleaved by the target protease and the cleavage is read out electrochemically. Two questions decide whether a candidate peptide is useful: \emph{is it cleaved by the target}, and \emph{is it cleaved by the other proteases likely to be present}? This document describes the software tool that answers both questions from sequence alone, how it was built, and how far its answers can be trusted.

\paragraph{What the tool does.} Given an eight-residue window, the tool returns for each protease in a chosen panel (i) a cleavage score, (ii) that score expressed as a percentile of random background sequences, and (iii) a specificity margin: the target's percentile minus the highest competitor percentile. It can also rank a pool of candidates, scan a protein sequence, and score every single-residue substitution of a peptide. A browser version with three-dimensional structure views accompanies this report.

\begin{keybox}
\begin{itemize}
\item \textbf{Seven trainable proteases.} MMP-2, MMP-8, MMP-9, MMP-14, neutrophil elastase (NE), LasB, and a group of digestive serine proteases; $10{,}472$ confirmed positive windows in total.
\item \textbf{Discrimination is real and leak-checked.} Cluster-aware cross-validated AUROC is \aucConsMMPnine{} for MMP-9 and \aucConsNE{} for NE (Table~\ref{tab:main}); random and cluster-aware folds differ by at most \maxLeakGap{} AUROC, and label-shuffled controls give \shuffleMin--\shuffleMax{}.
\item \textbf{The physicochemical encoding is not always best.} On data-rich proteases a tuning-free PSSM and a one-hot model beat the $64$-feature model by $0.02$--$0.03$ AUROC; on the small sets (NE, LasB) the physicochemical model is significantly better than the one-hot logistic regression. The shipped score therefore combines both (Section~\ref{sec:res-models}).
\item \textbf{Training against other proteases matters.} It raises cross-protease discrimination (for example LasB from $0.80$ to $0.86$ AUROC) at a cost of about $0.01$--$0.02$ against random background (Section~\ref{sec:res-neg}).
\item \textbf{The model recovers known biology it was never shown.} It ranks Pro at P3 first and Leu/Ile at P1$'$ first and second for MMP-9, matching the phage-display consensus of Kridel et al.\ (2001), and reproduces the direction of two designed mutants (Section~\ref{sec:bio-val}).
\item \textbf{Selectivity among MMPs is weak.} MMP-9 versus MMP-2 discrimination from a P4--P4$'$ window is only modestly above chance ($0.59$ AUROC in one direction, $0.68$ in the other; Table~\ref{tab:selmat}); this is a property of the problem and the data, not only of the model (Section~\ref{sec:contam}).
\item \textbf{A design flaw was found and fixed.} Three of the four candidate sensor cores were fed to the model shifted by one residue; with correct alignment the design-analysis condition number moves from $61.7$ to $16.1$ (Section~\ref{sec:design}).
\end{itemize}
\end{keybox}

\paragraph{Scope.} The models see only the eight residues around the scissile bond. They do not model the surrounding construct, three-dimensional structure, enzyme concentration, buffer conditions, or kinetics. A score is a statement about sequence resemblance to observed cleavage sites, not a rate constant (Section~\ref{sec:limits}).
